% ═══════════════════════════════════════════════════════════════
% SPANISCH MUSTERLÖSUNG: Verbos regulares -ER e -IR
% ═══════════════════════════════════════════════════════════════
% Thema: Regelmäßige Verben der 2. und 3. Konjugation
% Klasse 7 | Schulkontext | 90 Minuten
% ═══════════════════════════════════════════════════════════════

\documentclass[10pt, a4paper]{article}

\usepackage[utf8]{inputenc}
\usepackage[T1]{fontenc}
\usepackage[ngerman]{babel}

% --- SCHRIFTEN ---
\usepackage{palatino}
\usepackage{helvet}
\newcommand{\sanstext}[1]{{\fontfamily{phv}\selectfont #1}}

\usepackage{microtype}
\usepackage[margin=1.4cm, top=1.6cm, bottom=1.3cm]{geometry}
\usepackage{enumitem}
\usepackage{tabularx}
\usepackage{booktabs}
\usepackage{array}
\usepackage{multicol}
\usepackage{tikz}
\usetikzlibrary{calc}

\usepackage{fancyhdr}
\usepackage{lastpage}
\usepackage{needspace}

\usepackage{xcolor}
\usepackage{amssymb}
\usepackage{tcolorbox}
\tcbuselibrary{skins}
\usepackage{ulem}  % für Durchstreichen

% ═══════════════════════════════════════════════════════════════
% FARBEN
% ═══════════════════════════════════════════════════════════════

\definecolor{kopfzeile}{HTML}{1A1A1A}
\definecolor{akzent}{HTML}{C41E3A}          % Spanisch-Karminrot
\definecolor{hellgrau}{HTML}{F2F2F2}
\definecolor{mittelgrau}{HTML}{777777}
\definecolor{dunkelgrau}{HTML}{444444}
\definecolor{endungER}{HTML}{1565C0}        % Blau für -ER
\definecolor{endungIR}{HTML}{6A1B9A}        % Violett für -IR
\definecolor{endungAR}{HTML}{2E7D32}        % Grün für -AR
\definecolor{loesung}{HTML}{C41E3A}         % Rot für Lösungen

% ═══════════════════════════════════════════════════════════════
% VARIABLEN
% ═══════════════════════════════════════════════════════════════

\newcommand{\thema}{Verbos regulares: -ER e -IR}
\newcommand{\klasse}{7}
\newcommand{\maxpunkte}{36}

% ═══════════════════════════════════════════════════════════════
% KOPF- UND FUSSZEILE
% ═══════════════════════════════════════════════════════════════

\pagestyle{fancy}
\fancyhf{}
\renewcommand{\headrulewidth}{0pt}
\renewcommand{\footrulewidth}{0pt}
\cfoot{\sanstext{\footnotesize\textcolor{mittelgrau}{\thepage\,/\,\pageref{LastPage}}}}

% ═══════════════════════════════════════════════════════════════
% BOXEN
% ═══════════════════════════════════════════════════════════════

% Grammatik-Box
\newtcolorbox{grammatikbox}{
    colback=hellgrau,
    colframe=hellgrau,
    boxrule=0pt,
    arc=0pt,
    left=8pt, right=8pt, top=8pt, bottom=6pt,
    borderline north={1.5pt}{0pt}{dunkelgrau}
}

% Merk-Box
\newtcolorbox{merkbox}{
    colback=hellgrau,
    colframe=hellgrau,
    boxrule=0pt,
    arc=0pt,
    left=8pt, right=8pt, top=8pt, bottom=6pt,
    borderline north={1.5pt}{0pt}{akzent}
}

% ═══════════════════════════════════════════════════════════════
% BEFEHLE
% ═══════════════════════════════════════════════════════════════

% Spanisch: Palatino fett
\newcommand{\es}[1]{\textbf{#1}}

% Deutsch: klein, kursiv, grau
\newcommand{\de}[1]{{\small\textit{\textcolor{mittelgrau}{#1}}}}

% Lösung: rot, fett
\newcommand{\lsg}[1]{\textcolor{loesung}{\textbf{#1}}}

% Lücke
\newcommand{\luecke}[1]{\rule{#1}{0.4pt}}

% Punktebox
\newcommand{\punktebox}[1]{\hfill\sanstext{\small[\lsg{#1}/#1]}}

% Aufgabentitel
\newcommand{\aufgabe}[2]{%
    \needspace{4\baselineskip}%
    \vspace{10pt}%
    \noindent\sanstext{\textbf{#1}\quad#2}%
    \vspace{5pt}%
    \nopagebreak[4]%
}

% Teilaufgaben
\newlist{teil}{enumerate}{1}
\setlist[teil]{label=\alph*), leftmargin=1.8em, itemsep=4pt, parsep=0pt, topsep=3pt}

% Endungen farbig
\newcommand{\ender}[1]{\textcolor{endungER}{\textbf{#1}}}
\newcommand{\endir}[1]{\textcolor{endungIR}{\textbf{#1}}}
\newcommand{\endar}[1]{\textcolor{endungAR}{\textbf{#1}}}

\setlength{\parindent}{0pt}
\setlength{\parskip}{4pt}
\setlength{\columnsep}{24pt}

% ═══════════════════════════════════════════════════════════════
% DOKUMENT
% ═══════════════════════════════════════════════════════════════

\begin{document}

% ═══════════════════════════════════════════════════════════════
% HEADER
% ═══════════════════════════════════════════════════════════════

\noindent
\begin{tikzpicture}[remember picture, overlay]
    \fill[kopfzeile] (current page.north west) rectangle +(\paperwidth, -1cm);
    \node[anchor=west, text=white, font=\sanstext\large\bfseries]
        at ([xshift=1.4cm, yshift=-0.5cm]current page.north west) {\thema};
    \node[anchor=east, text=white, font=\sanstext\small]
        at ([xshift=-1.4cm, yshift=-0.5cm]current page.north east)
        {Spanisch\,·\,Klasse \klasse\,·\,\textcolor{loesung}{Musterlösung}};
\end{tikzpicture}

\vspace{0.5cm}

% --- Name/Datum ---
\noindent
\sanstext{\small\textbf{Name:}} \lsg{Musterlösung} \hfill
\sanstext{\small\textbf{Datum:}} \luecke{2.5cm}

\vspace{8pt}

% ═══════════════════════════════════════════════════════════════
% GRAMMATIK-ÜBERSICHT
% ═══════════════════════════════════════════════════════════════

\begin{grammatikbox}
\begin{multicols}{2}

% --- Endungen -ER ---
{\large\bfseries Verbos en \ender{-ER}} \de{2. Konjugation}

\vspace{3pt}

{\small
\begin{tabular}{@{}ll@{}}
yo & \ender{-o} \\
tú & \ender{-es} \\
él/ella & \ender{-e} \\
\end{tabular}
\hfill
\begin{tabular}{@{}ll@{}}
nosotros & \ender{-emos} \\
vosotros & \ender{-éis} \\
ellos/ellas & \ender{-en} \\
\end{tabular}}

\vspace{4pt}

{\footnotesize\textit{Beispiele:} \es{leer, comprender}}

\columnbreak

% --- Endungen -IR ---
{\large\bfseries Verbos en \endir{-IR}} \de{3. Konjugation}

\vspace{3pt}

{\small
\begin{tabular}{@{}ll@{}}
yo & \endir{-o} \\
tú & \endir{-es} \\
él/ella & \endir{-e} \\
\end{tabular}
\hfill
\begin{tabular}{@{}ll@{}}
nosotros & \endir{-imos} \\
vosotros & \endir{-ís} \\
ellos/ellas & \endir{-en} \\
\end{tabular}}

\vspace{4pt}

{\footnotesize\textit{Beispiele:} \es{abrir, escribir}}

\end{multicols}
\end{grammatikbox}

\vspace{2pt}

% ═══════════════════════════════════════════════════════════════
% AUFGABEN MIT LÖSUNGEN
% ═══════════════════════════════════════════════════════════════

\begin{multicols}{2}

% --- AUFGABE 1 ---
\aufgabe{1}{Konjugiere \es{leer} \de{lesen}.} \punktebox{6}

\vspace{4pt}

{\small
\renewcommand{\arraystretch}{1.5}
\begin{tabular}{@{}l@{\hspace{8pt}}p{3.2cm}@{}}
yo & le\lsg{o} \\
tú & le\lsg{es} \\
él/ella & le\lsg{e} \\
nosotros & le\lsg{emos} \\
vosotros & le\lsg{éis} \\
ellos/ellas & le\lsg{en} \\
\end{tabular}}

\vspace{6pt}

% --- AUFGABE 2 ---
\aufgabe{2}{Konjugiere \es{escribir} \de{schreiben}.} \punktebox{6}

\vspace{4pt}

{\small
\renewcommand{\arraystretch}{1.5}
\begin{tabular}{@{}l@{\hspace{8pt}}p{3.2cm}@{}}
yo & escrib\lsg{o} \\
tú & escrib\lsg{es} \\
él/ella & escrib\lsg{e} \\
nosotros & escrib\lsg{imos} \\
vosotros & escrib\lsg{ís} \\
ellos/ellas & escrib\lsg{en} \\
\end{tabular}}

\vspace{6pt}

% --- AUFGABE 3 ---
\aufgabe{3}{Setze die richtige Verbform ein.} \punktebox{6}

{\small Benutze: \es{leer, comprender, abrir, escribir}}

\vspace{4pt}

\begin{teil}
    \item María \lsg{lee} un libro. \de{lesen}

    \vspace{2pt}

    \item Nosotros \lsg{abrimos} la ventana. \de{öffnen}

    \vspace{2pt}

    \item ¿Tú \lsg{comprendes} la tarea? \de{verstehen}

    \vspace{2pt}

    \item Los alumnos \lsg{escriben} en el cuaderno. \de{schreiben}

    \vspace{2pt}

    \item Yo \lsg{abro} el libro de español. \de{öffnen}

    \vspace{2pt}

    \item Vosotros \lsg{comprendéis} muy bien. \de{verstehen}
\end{teil}

\columnbreak

% --- AUFGABE 4 ---
\aufgabe{4}{Gemischte Übung: -AR, -ER, -IR} \punktebox{8}

{\small Ergänze die Verbformen. Achte auf die Konjugation!}

\vspace{3pt}

{\footnotesize
\fbox{\parbox{0.9\linewidth}{
\es{hablar} · \es{escuchar} · \es{mirar} · \es{explicar} · \es{practicar} \\
\es{leer} · \es{comprender} · \es{abrir} · \es{escribir}
}}}

\vspace{4pt}

\begin{teil}
    \item El profesor \lsg{explica} la gramática. \de{erklären}

    \vspace{2pt}

    \item Los alumnos \lsg{escuchan} al profesor. \de{zuhören}

    \vspace{2pt}

    \item Yo \lsg{hablo} español muy bien. \de{sprechen}

    \vspace{2pt}

    \item Nosotros \lsg{leemos} un texto. \de{lesen}

    \vspace{2pt}

    \item ¿Tú \lsg{miras} la pizarra? \de{anschauen}

    \vspace{2pt}

    \item Ellos \lsg{practican} los verbos. \de{üben}

    \vspace{2pt}

    \item María \lsg{escribe} una carta. \de{schreiben}

    \vspace{2pt}

    \item Vosotros \lsg{abrís} el libro. \de{öffnen}
\end{teil}

\vspace{6pt}

% --- AUFGABE 5 ---
\aufgabe{5}{Was passt nicht? Streiche durch.} \punktebox{4}

\vspace{4pt}

{\small
\begin{teil}
    \item \es{leo} -- \es{lees} -- \sout{\lsg{lea}} -- \es{leemos}

    \vspace{2pt}

    \item \es{escribo} -- \es{escribe} -- \sout{\lsg{escribas}} -- \es{escriben}

    \vspace{2pt}

    \item \es{abrimos} -- \es{abrís} -- \sout{\lsg{abréis}} -- \es{abren}

    \vspace{2pt}

    \item \es{comprendo} -- \es{comprendemos} -- \sout{\lsg{comprendimos}} -- \es{comprendéis}
\end{teil}}

\end{multicols}

% ═══════════════════════════════════════════════════════════════
% AUFGABE 6: Volle Breite
% ═══════════════════════════════════════════════════════════════

\vspace{2pt}
\noindent\textcolor{hellgrau}{\rule{\textwidth}{1.5pt}}
\vspace{2pt}

\aufgabe{6}{Schreibe 4 Sätze über deinen Schulalltag.} \punktebox{6}

{\small Benutze mindestens je 1 Verb auf -AR, -ER und -IR!}

\vspace{3pt}

{\footnotesize\textit{Ejemplo:} En clase, yo leo libros y escribo en el cuaderno. El profesor explica la gramática.}

\vspace{6pt}

{\small\lsg{Mögliche Lösungen (individuelle Antworten):}}

\vspace{4pt}

1. \lsg{En la clase de español, yo escucho al profesor.}

\vspace{8pt}

2. \lsg{María lee un libro y comprende el texto.}

\vspace{8pt}

3. \lsg{Nosotros abrimos los cuadernos y escribimos.}

\vspace{8pt}

4. \lsg{El profesor explica la gramática y los alumnos practican.}

% ═══════════════════════════════════════════════════════════════
% MERKKASTEN
% ═══════════════════════════════════════════════════════════════

\vspace{8pt}

\begin{merkbox}
\sanstext{\textbf{Recuerda – Unterschied -ER / -IR}}

\vspace{3pt}

\begin{multicols}{2}

{\small
\renewcommand{\arraystretch}{1.2}
\begin{tabular}{@{}lcc@{}}
& \ender{-ER} & \endir{-IR} \\
\hline
nosotros & \ender{-emos} & \endir{-imos} \\
vosotros & \ender{-éis} & \endir{-ís} \\
\end{tabular}}

\columnbreak

{\small
Die anderen Formen sind \textbf{identisch}:\\
-o, -es, -e, -en\\[2pt]
{\footnotesize\textbf{Tipp:} Merke dir nur nosotros/vosotros!}}

\end{multicols}
\end{merkbox}

% ═══════════════════════════════════════════════════════════════
% FOOTER
% ═══════════════════════════════════════════════════════════════

\vfill

\noindent\rule{\textwidth}{0.5pt}

\vspace{4pt}

\hfill\sanstext{\textbf{Punkte:}} \lsg{36} / \maxpunkte

\end{document}
